\documentclass[aps,prl,twocolumn,10pt,superscriptaddress]{revtex4-2}

\usepackage{labreport}

\begin{document}

% Define notebook cells extraction macros (if using auto-extracted code)
% % auto_nb_code.tex — generated by nb_to_report.py
% Place % auto_nb_code.tex — generated by nb_to_report.py
% Place % auto_nb_code.tex — generated by nb_to_report.py
% Place \input{auto_nb_code} inside Analysis Code section.

\paragraph{Setup}
\begin{Code}[H]
  \lstinputlisting[language=Python]{pycode/nb_cells/_nb_setup.py}
  \caption*{Notebook cell: Setup}
\end{Code}

\paragraph{Computation}
\begin{Code}[H]
  \lstinputlisting[language=Python]{pycode/nb_cells/_nb_computation.py}
  \caption*{Notebook cell: Computation}
\end{Code}
 inside Analysis Code section.

\paragraph{Setup}
\begin{Code}[H]
  \lstinputlisting[language=Python]{pycode/nb_cells/_nb_setup.py}
  \caption*{Notebook cell: Setup}
\end{Code}

\paragraph{Computation}
\begin{Code}[H]
  \lstinputlisting[language=Python]{pycode/nb_cells/_nb_computation.py}
  \caption*{Notebook cell: Computation}
\end{Code}
 inside Analysis Code section.

\paragraph{Setup}
\begin{Code}[H]
  \lstinputlisting[language=Python]{pycode/nb_cells/_nb_setup.py}
  \caption*{Notebook cell: Setup}
\end{Code}

\paragraph{Computation}
\begin{Code}[H]
  \lstinputlisting[language=Python]{pycode/nb_cells/_nb_computation.py}
  \caption*{Notebook cell: Computation}
\end{Code}


\title{QNT SCI 410 - Lab N: [Title]}

\author{Gregory Sinaga}
\email{gregorysinaga@ucla.edu}
\affiliation{Department of Physics and Astronomy, University of California, Los Angeles}

\author{[Lab Partner]}
\affiliation{Department of Physics and Astronomy, University of California, Los Angeles}

\date{\today}

\begin{abstract}
% Include key quantitative results: measured values with uncertainties,
% statistical significance of violations, and figure-of-merit comparisons.
% Target: 150-250 words. Every number should appear here.
[Abstract text with key numerical results and uncertainties.]
\end{abstract}

\maketitle

% ===========================================================================
\section{Introduction and Background}
% ===========================================================================

\begin{guidance}
Provide context for the experiment: what physical phenomena are being explored,
how this connects to previous labs, and what the key theoretical predictions are.
Use \emph{} for introduced terms. Cite lecture notes and textbooks.
Target: 1-2 columns. End with a roadmap paragraph listing the experiments.
\end{guidance}

% [Introduction text]

% Subsections for key concepts:
% \subsection*{Concept A}
% \subsection*{Concept B}

% ===========================================================================
\section{Theory}
\label{sec:theory}
% ===========================================================================

\begin{guidance}
Formal derivations and equations needed to interpret the data.
Number all key equations with \verb|\begin{equation}...\end{equation}|.
Cross-reference equations in the text with \verb|\autoref{eq:label}|.
Use notation consistent with lecture notes (Campbell) and manuals (qutools).
\end{guidance}

% \subsection{Subsection A}
% \subsection{Subsection B}

% ===========================================================================
\section{Experimental Procedure}
% ===========================================================================

\begin{guidance}
Split into subsections per experiment/lab session.
Each subsection needs: Equipment table, detailed procedure (numbered steps),
and a TikZ or photograph figure of the setup.
Use \verb|\bomtable| or inline equipment tables.
\end{guidance}

% \subsection{Lab Na --- [Experiment A]}
% \label{subsec:proc-a}

% \noindent\textbf{Equipment:}
% \begin{table*}[t]
% ... equipment table ...
% \end{table*}

% \paragraph{Detailed Procedure:}
% \begin{enumerate}
%   \item \textbf{Step 1:} ...
%   \item \textbf{Step 2:} ...
% \end{enumerate}

% ===========================================================================
\section{Data Acquisition, Analysis and Discussion}
% ===========================================================================

\begin{guidance}
Present results with discussion interleaved (not separated).
For each experiment:
  1. State the experimental goal in one sentence (\emph{Experimental Goal.})
  2. Present raw data (table or figure)
  3. Show analysis (fits, derived quantities with uncertainties)
  4. Discuss: what does this mean physically? Compare to theory.
  5. Identify systematic effects and limitations.

Every figure must be referenced in text. Every number needs uncertainty.
Every quantitative claim needs physical interpretation (2-3 sentences).
Use \verb|\csvdatatable| for data tables from CSV files.
Use \verb|\pythonsnippet| for analysis code listings.
\end{guidance}

% \subsection{Lab Na --- [Experiment A Results]}
% \subsubsection{Raw Data}
% \subsubsection{Analysis}
% \subsubsection{Discussion}

% \subsection{Uncertainty Analysis and Error Propagation}
% \subsubsection{Statistical Sources}
% \subsubsection{Systematic Uncertainties}

% \subsection{Goodness-of-Fit Assessment}
% chi-squared analysis

% ===========================================================================
\section{Conclusion}
% ===========================================================================

\begin{guidance}
Summarize key findings from each experiment with numbers.
State what was demonstrated and its significance.
Acknowledge limitations (loopholes, systematic effects).
Suggest improvements for future work.
Target: 1 column.
\end{guidance}

% ===========================================================================
\appendix
\onecolumngrid
\section{Raw Data Tables}
\label{app:raw-data}
% ===========================================================================

% Full raw data tables go here (if too large for main body).

% ===========================================================================
\bibliographystyle{apsrev4-2}
\bibliography{refs}
% ===========================================================================

\end{document}
